
\section{Threats to validity}

We assess potential threats to the validity of our work, considering the construct, internal, conclusion, and external factors as outlined in \cite{wohlin2012experimentation}. The \textbf{construct threat} primarily concerns hyperparameters settings and evaluation metrics. To address this, we conduct a sensitivity analysis for all methods that require specific parameters and use established evaluation metrics. The \textbf{internal threat} concerns the framework implementation, where bugs may affects the results reliability. To mitigate, we use established Python packages, rigorous testing, and repeat the experiments multiple times. The \textbf{conclusion threat} is tied to the fault types as microservices can experience various faults \cite{mariani2018localizing}. We use four common faults to evaluate and demonstrate BARO's superior performance. Furthermore, our framework relies on a set of assumptions discussed and validated in previous works on microservice systems, as described in Sec \ref{sec:assumption-causal}. If a system meets these assumptions, BARO could be applied. Expanding BARO to work with other types of systems, e.g., distributed database systems, could be a potential future work. The \textbf{external threat} is related to the deployment of microservice applications and data collection strategies. In this paper, we employ a single deployment setting for each microservice system and a single data collection strategy, which potentially limits the generality of our work.  However, we deploy the benchmark systems on the real 5-node Kubernetes clusters and repeat the experiments multiple times to get comprehensive datasets. In addition, we use popular benchmark microservice applications, such as Online Boutique, Sock Shop, and Train Ticket, which are widely recognized in academia for testing microservices-related methods~\cite{Jinjin2018Microscope, Azam2022rcd, wu2021microdiag, Xin2023CausalRCA, wu2022automatic, he2022graph, dan2021practical, yu2021microrank, zhou2018trainticket, Wang2021evalcausal}. Furthermore, we adopt Prometheus, an open-source tool for real-time monitoring, and collect service-level and resource-level metrics that present the status of a running microservice application. This choice helps mitigate this threat.


